\section{Giá trị lượng giác của một góc}

\subsection{Đề 1}
\Opensolutionfile{ans}[ans/ans-0-B1]
\caulc
\begin{ex}%[0H4N1-1]
	Cho $0^{\circ}  < \alpha < 90^{\circ} $. Khẳng định nào sau đây đúng?
	\choice
	{$\cot \left(90^{\circ} -\alpha \right)=-\tan \alpha$}
	{\True $\cos \left(90^{\circ} -\alpha \right)=\sin \alpha$}
	{$\sin \left(90^{\circ} -\alpha \right)=-\cos \alpha$}
	{$\tan \left(90^{\circ} -\alpha \right)=-\cot \alpha$}
	\loigiai{
		Vì $\alpha$ và $\left(90^{\circ} -\alpha \right)$ là hai cung phụ nhau.
	}
\end{ex}%[0H4N1-2]
%%%==============EX_2============%%%
\begin{ex}
	Tính giá trị biểu thức $P=\sin 30^{\circ} \cos 60^{\circ}+\sin 60^{\circ} \cos 30^{\circ}$
	\choice
	{\True $P=1$}
	{$P=0$}
	{$P=\sqrt{3}$}
	{$P=-\sqrt{3}$}
	\loigiai{
		Ta có $P=\sin 30^{\circ} \cdot \sin 30^{\circ}+\cos 30^{\circ} \cdot \cos 30^{\circ}=\sin^2 30^{\circ}+\cos^2 30^{\circ}=1$.
	}
\end{ex}
%%%==============EX_3============%%%
\begin{ex}%[0H4H1-1]
	Biết $A$, $B$, $C$ là ba góc của tam giác $ABC$, mệnh đề nào sau đây là đúng?
	\choice
	{$\cos \left(A+C\right)=\cos B$}
	{\True $\tan \left(A+C\right)=-\tan B$}
	{$\cot \left(A+C\right)=\cot B$}
	{$\sin \left(A+C\right)=-\sin B$}
	\loigiai{
		Xét phương án A: $\cos \left(A+C\right)=\cos \left(\pi-B\right)=-\cos B \Rightarrow$ A sai.\\
		Xét phương án B: $\tan \left(A+C\right)=\tan \left(\pi-B\right)=-\tan B$. Vậy B đúng.\\
		Xét phương án C: $\cot \left(A+C\right)=\cot \left(\pi-B\right)=-\cot B \Rightarrow$ C sai.\\
		Xét phương án D: $\sin \left(A+C\right)=\sin \left(\pi-B\right)=\sin B \Rightarrow$ D sai.
	}
\end{ex}


\begin{ex}%[0H4N1-1]
	Cho $\alpha$ là góc tù. Điều khẳng định nào sau đây là đúng?
	\choice
	{$\sin \alpha < 0$}
	{$\cos \alpha > 0$}
	{\True $\tan \alpha < 0$}
	{$\cot \alpha > 0$}
	\loigiai{
		Góc tù có điểm biểu diễn thuộc góc phần tư thứ II, có giá trị $\sin \alpha > 0$, còn $\cos \alpha$, $\tan \alpha$ và $\cot \alpha$ đều nhỏ hơn $0$.
	}
\end{ex}

\begin{ex}%[0H4N1-1]
	Trong các đẳng thức sau, đẳng thức nào đúng?
	\choice
	{$\sin \left(180^{\circ}-\alpha \right)=\cos \alpha$}
	{$\sin \left(180^{\circ}-\alpha \right)=-\cos \alpha$}
	{$\sin \left(180^{\circ}-\alpha \right)=-\sin \alpha$}
	{\True $\sin \left(180^{\circ}-\alpha \right)=\sin \alpha$}
	\loigiai{
	}
\end{ex}


\begin{ex}%[0H4V1-2]
	Tam giác đều $ABC$ có đường cao $AH$. Khẳng định nào sau đây là đúng?
	\choice
	{$\cos \widehat{BAH}=\dfrac{1}{\sqrt{3}}$}
	{\True $\sin \widehat{ABC}=\dfrac{\sqrt{3}}{2}$}
	{$\sin \widehat{AHC}=\dfrac{1}{2}$}
	{$\sin \widehat{BAH}=\dfrac{\sqrt{3}}{2}$}
	\loigiai{
		Tam giác $ABC$ là tam giác đều nên có các góc bằng $60^{\circ}$ nên dễ thấy C đúng vì
		\[\sin \widehat{ABC}=\sin 60^{\circ}=\dfrac{\sqrt{3}}{2}.
		\]
	}
\end{ex}


\begin{ex}%[0H4H1-2]
	Tam giác $ABC$ vuông ở $A$ có góc $\widehat{B}=30^{\circ}$. Khẳng định nào sau đây là sai?
	\choice
	{\True $\cos B=\dfrac{1}{\sqrt{3}}$}
	{$\sin C=\dfrac{\sqrt{3}}{2}$}
	{$\cos C=\dfrac{1}{2}$}
	{$\sin B=\dfrac{1}{2}$}
	\loigiai{
		Dễ thấy A sai do $\cos B=\cos 30^{\circ}=\dfrac{\sqrt{3}}{2}$.
	}
\end{ex}

\begin{ex}%[0H4N1-2]
	Trong các đẳng thức sau đây đẳng thức nào là đúng?
	\choice
	{$\sin 150^{\circ}=-\dfrac{\sqrt{3}}{2}$}
	{$\cos 150^{\circ}=\dfrac{\sqrt{3}}{2}$}
	{\True $\tan 150^{\circ}=-\dfrac{1}{\sqrt{3}}$}
	{$\cot 150^{\circ}=\sqrt{3}$}
	\loigiai{
		Dựa vào giá trị lượng giác của các cung bù nhau. Dễ thấy phương án đúng là}
	{\True Ta có $\sin 150^{\circ}=\sin 30^{\circ}=\dfrac{1}{2}$,
		
		$\cos 150^{\circ}=-\cos 30^{\circ}=-\dfrac{\sqrt{3}}{2}$,
		
		
		$\tan 150^{\circ}=-\tan 30^{\circ}=-\dfrac{1}{\sqrt{3}}$ và $\cot 150^{\circ}=-\cot 30^{\circ}=-\sqrt{3}$.
	}
\end{ex}


\begin{ex}%[0H4H1-2]
	Cho góc $\alpha \in \left(90^{\circ};180^{\circ} \right)$. Khẳng định nào sau đây đúng?
	\choice
	{$\sin \alpha$ và $\cot \alpha$ cùng dấu}
	{\True Tích $\sin \alpha \cdot \cot \alpha$ mang dấu âm}
	{Tích $\sin \alpha  \cdot \cos \alpha$ mang dấu dương}
	{$\sin \alpha$ và $\tan \alpha$ cùng dấu}
	\loigiai{
		Với $\alpha \in \left(90^{\circ};180^{\circ} \right)$, ta có $\sin \alpha > 0,\cos \alpha < 0$ suy ra: $\tan \alpha < 0,\cot \alpha < 0$
		Vậy $\sin \alpha \cdot \cot \alpha < 0$
	}
\end{ex}

\begin{ex}%[0H4N1-2]
	Hai góc nhọn $\alpha$ và $\beta$ phụ nhau, hệ thức nào sau đây là sai?
	\choice
	{$\sin \alpha=\cos \beta$}
	{$\tan \alpha=\cot \beta$}
	{$\cot \beta=\dfrac{1}{\cot \alpha}$}
	{\True $\cos \alpha=-\sin \beta$}
	\loigiai{
		\[\cos \alpha=\cos \left(90^{^{\circ}}-\beta \right)=\sin \beta.
		\]
	}
\end{ex}


\begin{ex}%[0H4V1-2]
	Cho $\sin \alpha=\dfrac{1}{3}$, với $90^{\circ} < \alpha < 180^{\circ}$. Tính $\cos \alpha$.
	\choice
	{$\cos \alpha=\dfrac{2}{3}$}
	{$\cos \alpha=-\dfrac{2}{3}$}
	{$\cos \alpha=\dfrac{2\sqrt{2}}{3}$}
	{\True $\cos \alpha=-\dfrac{2\sqrt{2}}{3}$}
	\loigiai{
		Ta có $\cos^2 \alpha=1-\sin^2 \alpha$ $=1-\left(\dfrac{1}{3} \right)^2=\dfrac{8}{9}$.
		Mặt khác $90^{\circ} < \alpha < 180^{\circ}$ nên $\cos \alpha=-\dfrac{2\sqrt{2}}{3}$.
	}
\end{ex}


\begin{ex}%[0H4V1-2]
	Cho biết $\cos \alpha=-\dfrac{2}{3}$. Tính $\tan \alpha$?
	\choice
	{$\dfrac{5}{4}$}
	{$-\dfrac{5}{2}$}
	{$\dfrac{\sqrt{5}}{2}$}
	{\True $-\dfrac{\sqrt{5}}{2}$}
	\loigiai{
		Do $\cos \alpha < 0\Rightarrow \tan \alpha < 0$.
		Ta có: $1+\tan^2 \alpha=\dfrac{1}{\cos^2 \alpha} \Leftrightarrow \tan^2 \alpha=\dfrac{5}{4} \Rightarrow \tan \alpha=-\dfrac{\sqrt{5}}{2}$.
	}
\end{ex}
\Closesolutionfile{ans}
\indapan{6}{ans/ans-0-B1}


\cauds % câu đúng sai
\Opensolutionfile{ans}[ans/ans-0-B1-DS]

\begin{ex}%[0H4V1-3]
	Xét tính \textbf{đúng}, \textbf{sai} của các đẳng thức sau:
	\choiceTF
	{\True $\left(\sin^2 \alpha+\cos^2 \alpha \right)^2=1$}
	{\True $1+2\sin \alpha \cdot \sin \alpha=(\sin \alpha+\cos \alpha)^2$}
	{\True $1-2\sin \alpha \cdot \sin \alpha=(\sin \alpha-\cos \alpha)^2$}
	{\True $1-2\sin^2 \alpha \cdot \sin^2 \alpha=\sin \alpha^4+\cos^4 \alpha$}
	\loigiai{
		\begin{itemchoice}
			\itemch Do $\sin^2 \alpha+\cos^2 \alpha=1$ nên $\left(\sin^2 \alpha+\cos^2 \alpha \right)^2=1$
			\itemch 	$1+2\sin \alpha \cdot \cos \alpha=\sin^2 \alpha+\cos^2 \alpha+2\sin \alpha \cdot \cos \alpha=(\sin \alpha+\cos \alpha)^2$
			\itemch 	$1-2\sin \alpha \cdot \cos \alpha=\sin^2 \alpha+\cos^2 \alpha-2\sin \alpha \cdot \cos \alpha=(\sin \alpha-\cos \alpha)^2$
			\itemch 		Ta có: $\sin^4 \alpha+\cos^4 \alpha+2\sin^2 \alpha \cdot \cos^2 \alpha=\left(\sin^2 \alpha+\cos^2 \alpha \right)^2=1$.
			
			Do đó: $1-2\sin^2 \alpha \cos^2 \alpha=\sin^4 \alpha+\cos^4 \alpha$.
		\end{itemchoice}
	}
\end{ex}

\begin{ex}%[0H4H1-2]
	Xét tính \textbf{đúng},\textbf{ sai} của các đẳng thức sau
	\choiceTF
	{\True $A=4\sin 30^{^{\circ}}+(\sqrt{3})^3 \cdot \tan 30^{^{\circ}}=5$}
	{$B=\dfrac{1}{\sqrt{3}} \cos 30^{^{\circ}}-3\sqrt{2} \sin 45^{^{\circ}}+\cot 45^{^{\circ}}=\dfrac{3}{2}$}
	{\True$C=\sin^2 60^{^{\circ}}+\tan^2 30^{^{\circ}}-2=-\dfrac{11}{12}$}
	{$D=2\left(\sin 150^{^{\circ}}+\sqrt{3} \right)+\dfrac{1}{2\sqrt{2}} \cos 135^{^{\circ}}-3\tan 150^{^{\circ}}=\dfrac{3}{4}+2\sqrt{3}$}
	\loigiai{
		\begin{itemchoice}
			\itemch $A=4\cdot\dfrac{1}{2}+3\sqrt{3}\cdot \dfrac{\sqrt{3}}{3}=5$
			\itemch $B=\dfrac{1}{\sqrt{3}}\cdot \dfrac{\sqrt{3}}{2}-3\sqrt{2}\cdot \dfrac{\sqrt{2}}{2}+1=-\dfrac{3}{2}$
			\itemch $C=\left(\dfrac{\sqrt{3}}{2}\right)^2+\left( \dfrac{\sqrt{3}}{3}\right)^2-2=-\dfrac{11}{12}$
			\itemch $2\left(\dfrac{1}{2}+\sqrt{3}\right)+\dfrac{1}{2\sqrt{2}}\cdot\left(-\dfrac{\sqrt{2}}{2}\right)-3\cdot\left(-\dfrac{\sqrt{3}}{3}\right)=\dfrac{3+12\sqrt{3}}{4}$
		\end{itemchoice}
	}
	
\end{ex}
%%%==============EX_4============%%%
\begin{ex}%[0H4V1-3]
	Cho $\sin \alpha=\dfrac{1}{3}$ với $90^{^{\circ}} < \alpha < 180^{^{\circ}}$. Xét tính \textbf{đúng}, \textbf{sai} các ý sau:
	\choiceTF
	{$\cos \alpha > 0$}
	{\True$\cos \alpha=-\dfrac{2\sqrt{2}}{3}$}
	{\True$\tan \alpha=-\dfrac{1}{2\sqrt{2}}$}
	{$\cot \alpha=2\sqrt{2}$}
	\loigiai{
		\begin{itemchoice}
			\itemch 	Vì $90^{^{\circ}} < \alpha < 180^{^{\circ}}$ nên $\cos \alpha < 0$
			\itemch 	Ta có $\sin^2 \alpha+\cos^2 \alpha=1\Rightarrow \cos^2 \alpha=1-\sin^2 \alpha=1-\left(\dfrac{1}{3} \right)^2=\dfrac{8}{9} \Rightarrow \cos \alpha=-\dfrac{2\sqrt{2}}{3}$
			\itemch 	Ta có $\tan \alpha=\dfrac{\sin \alpha}{\cos \alpha}=\dfrac{\dfrac{1}{3}}{-\dfrac{2\sqrt{2}}{3}}=-\dfrac{1}{2\sqrt{2}}$
			\itemch Ta có  $\cot \alpha = \dfrac{1}{\tan \alpha}=-2\sqrt{2}$
		\end{itemchoice}
		
		
		
		
	}
\end{ex}


\begin{ex}%[0H4V1-3]
	Cho $\tan \alpha=-\dfrac{3}{4}$, $90^{\circ}< \alpha < 180^{\circ}$.   Xét tính \textbf{đúng}, \textbf{sai} các ý sau:
	\choiceTF
	{$\cos \alpha > 0$}
	{\True$\cos \alpha = -\dfrac{4}{5}$}
	{\True$\cot \alpha=-\dfrac{4}{3}$}
	{$\sin \alpha=\dfrac{3}{5}$}
	\loigiai{
		\begin{itemchoice}
			\itemch Theo giả thiết $90^{\circ}<\alpha < 180^{\circ} \Rightarrow \cos \alpha <0$
			\itemch Ta có $\dfrac{1}{\cos^2}=1+\tan^2\alpha = 1 +\left( \dfrac{-3}{4}\right)^2=\dfrac{25}{16} \Rightarrow \cos^2 \alpha = \dfrac{16}{25}\Rightarrow \cos \alpha =-\dfrac{4}{5}$
			\itemch Ta có $\tan \alpha \cdot \cot \alpha = 1 \Rightarrow \cot \alpha = \dfrac{1}{\tan \alpha}=-\dfrac{4}{3} $
			\itemch Ta có $\tan \alpha = \dfrac{\sin \alpha}{\cos \alpha}\Rightarrow \sin \alpha = \tan \alpha \cdot \cos \alpha =-\dfrac{3}{4}\cdot -\dfrac{4}{5}=\dfrac{3}{5}$ 
		\end{itemchoice}
	}
\end{ex}		
\Closesolutionfile{ans}
\indapan{3}{ans/ans-0-B1-DS}

\Opensolutionfile{ans}[ans/ans-0-B1-KQ]
\caukq

\begin{ex}%[0H4C1-3]
	Cho $\cot \alpha=-3$. Tính giá trị biểu thức $P=\dfrac{\sin^3 \alpha+\cos^3 \alpha}{\sin \alpha-\cos \alpha}$. (Kết quả làm tròn đến hàng phần mười)
	\shortans{$-0{,}7$}
	\loigiai{
		Vì $\cot \alpha=-3$ nên $\sin \alpha \ne 0$. Chia cả tử và mẫu của $P$ cho $\sin^3 \alpha$, ta có:
		\[P=\dfrac{1+\cot^3 \alpha}{\dfrac{1}{\sin^2 \alpha} \left(1-\cot \alpha \right)}=\dfrac{1+\cot^3 \alpha}{\left(1+\cot^2 \alpha \right)\left(1-\cot \alpha \right)}=\dfrac{1+\left(-3\right)^3}{\left[1+\left(-3\right)^2 \right]\left[1-\left(-3\right)\right]}=-\dfrac{13}{20}
		\]
	}
\end{ex}

\begin{ex}%[0H4V1-3]
	Cho $\cos x=\dfrac{1}{2}$. Tính giá trị biểu thức $P=3\sin^2 x+4\cos^2 x$?
	\shortans{$3{,}25$}
	\loigiai{
		Ta có: $P=3\sin^2 x+4\cos^2 x=3\left(1-\cos^2 x\right)+4\cos^2 x=3+\cos^2 x=3+\dfrac{1}{4}=\dfrac{13}{4}$.
	}
\end{ex}

\begin{ex}%[0H4V1-3]
	Cho $\tan \alpha=\dfrac{1}{3}$. Tính giá trị của biểu thức $A=\dfrac{3\sin \alpha+4\cos \alpha}{2\sin \alpha-5\cos \alpha}$ thu được kết quả dạng $-\dfrac{a}{b}$ với $\dfrac{a}{b}$ là phân số tối giản. Tính $a+b$.
	\shortans{$28$}
	\loigiai{
		Vì $\tan \alpha=\dfrac{\sin \alpha}{\cos \alpha}=\dfrac{1}{3}$ nên $\cos \alpha \ne 0$.
		
		Chia cả tử và mẫu của $P$ cho $\cos \alpha$, ta được: $A=\dfrac{3\dfrac{\sin \alpha}{\cos \alpha}+4}{2\dfrac{\sin \alpha}{\cos \alpha}-5}=\dfrac{3\tan \alpha+4}{2\tan \alpha-5}=\dfrac{3\cdot \dfrac{1}{3}+4}{2\cdot \dfrac{1}{3}-5}=-\dfrac{15}{13}$. Vậy $a=15,b=13$ nên $a+b=28$.
	}
\end{ex}

\begin{ex}%[0H4V1-3]
	Tính giá trị biểu thức sau: $B=\tan 1^{^{\circ}} \tan 2^{^{\circ}} \tan 3^{^{\circ}} \ldots \tan 89^{^{\circ}}$.
	\shortans{$1$}
	\loigiai{
		$$
		\begin{aligned}
			B&=\tan 1^{^{\circ}} \tan 2^{^{\circ}} \tan 3^{^{\circ}} \ldots \tan 89^{^{\circ}}\\
			&=\left(\tan 1^{^{\circ}} \tan 89^{^{\circ}} \right)\cdot \left(\tan 2^{^{\circ}} \tan 88^{^{\circ}} \right)\ldots \left(\tan 44^{^{\circ}} \tan 46^{^{\circ}} \right)\cdot \tan 45^{^{\circ}}\\
			&=\left(\tan 1^{^{\circ}} \cdot \cot 1^{^{\circ}} \right)\cdot \left(\tan 2^{^{\circ}} \cot 2^{^{\circ}} \right)\ldots \left(\tan 44^{^{\circ}} \cot 44^{^{\circ}} \right)\cdot \tan 45^{^{\circ}}\\
			&=1\cdot1\ldots 1=1	
		\end{aligned}
		$$
		
	}
\end{ex}

\begin{ex}%[0H4C1-3]
	Tính giá trị biểu thức sau: $C=\sin^2 51^{^{\circ}}+\sin^2 55^{^{\circ}}+\sin^2 39^{^{\circ}}+\sin^2 35^{^{\circ}}$
	\shortans{$2$}
	\loigiai{
		$$\begin{aligned}
			C&=\left(\sin^2 51^{^{\circ}}+\sin^2 39^{^{\circ}} \right)+\left(\sin^2 55^{^{\circ}}+\sin^2 35^{^{\circ}} \right)\\
			&=\left[\sin^2 51^0+\sin^2 \left(90^{^{\circ}}-51^0 \right)\right]+\left[\sin^2 55^{^{\circ}}+\sin^2 \left(90^{^{\circ}}-55^{^{\circ}} \right)\right]\\
			&=\left(\sin^2 51^{^{\circ}}+\cos^2 51^{^{\circ}} \right)+\left(\sin^2 55^{^{\circ}}+\cos^2 55^{^{\circ}} \right)\\
			&=1+1=2
		\end{aligned}$$
	}
\end{ex}

\begin{ex}%[0H4N1-2]
	Tính giá trị của biểu thức: $A=\sin 60^{^{\circ}}+2\cos 30^{^{\circ}}-3\sin 45^{^{\circ}}$ thu được dạng $\dfrac{\sqrt{3}-a\sqrt{2}}{b}$ với $a,b\in \mathbb{N}$. Tính $a+b$.
	\shortans{ $5$}
	\loigiai{
		\[A=\dfrac{\sqrt{3}}{2}+2\cdot \dfrac{\sqrt{3}}{2}-3\cdot \dfrac{\sqrt{2}}{2}=\dfrac{\sqrt{3}-3\sqrt{2}}{2}.
		\]
		Vậy $a=3, b=2$ nên $a+b=5$.
	}
\end{ex}
\Closesolutionfile{ans}
\indapan{6}{ans/ans-0-B1-KQ}
\newpage 
\subsection{Đề 2}
\Opensolutionfile{ans}[ans/ans-0-B1]
\caulc
\begin{ex}%[0D3N1-2] %1
	Cho $0^\circ < \alpha < 90^\circ$. Khẳng định nào sau đây là đúng?
	\choice
	{ $\cot(90^\circ-\alpha)=-\tan\alpha$}
	{\True $\cos(90^\circ -\alpha = \sin \alpha)$}
	{$\sin(90^\circ - \alpha) =-\cos\alpha$}
	{$\tan(90^\circ-\alpha)= - \cot \alpha$}
	\loigiai{
		Vì $\alpha$ và $90^\circ - \alpha$ là hai cung phụ nhau nên theo tính chất giá trị lượng giác của hai cung phụ nhau ta có đáp án. }
\end{ex}	
\begin{ex}%[0D3N1-2] %2
	Tính giá trị biểu thức  $P = \sin 30^\circ \cos 60^\circ  + \sin 60^\circ \cos 30^\circ $.
	\choice
	{ \True $P=1$}
	{ $P=0$}
	{$P=\sqrt{3}$}
	{$P=-\sqrt{3}$}
	\loigiai{
		Ta có $P = \sin 30^\circ .\sin 30^\circ  + \cos 30^\circ .\cos30^\circ  = {\sin ^2}30^\circ  + {\cos ^2}30^\circ  = 1$. }
\end{ex}	
\begin{ex}%[0D3N1-2] %3
	Cho $\sin x + \cos x = m$. Tính theo $m$ giá trị của $M=\sin x . \cos x$.
	\choice
	{ $m^2 -1$}
	{\True $\dfrac{m^2 -1}{2}$}
	{$\dfrac{m^2 +1}{2}$}
	{$m^2 +1$}
	\loigiai{
		$\begin{gathered}
			\sin x + \cos x = m\,\, \Rightarrow \,{\left( {\sin x + \cos x} \right)^2}\, = {m^2} \hfill \\
			\Leftrightarrow \,\left( {{{\sin }^2}x + {{\cos }^2}x} \right) + 2\sin x.\cos x = {m^2}\, \hfill \\
			\Leftrightarrow \,1 + 2\sin x.\cos x = {m^2} \Leftrightarrow \,\sin x.\cos x = \dfrac{{{m^2} - 1}}{2} \hfill \\ 
		\end{gathered}. $\\
		Vậy $M = \dfrac{{{m^2} - 1}}{2}$. }
\end{ex}
\begin{ex}%[0D3N1-2] %4
	Biểu thức $P = {\cos ^4}\alpha  + {\cos ^2}\alpha {\sin ^2}\alpha  + {\sin ^2}\alpha $ có giá trị bằng
	\choice
	{\True $1$}
	{ $2$}
	{$-2$}
	{$-1$}
	\loigiai{
		$P = {\cos ^2}\alpha \left( {{{\cos }^2}\alpha  + {{\sin }^2}\alpha } \right) + {\sin ^2}\alpha  = {\cos ^2}\alpha  + {\sin ^2}\alpha  = 1$. }
\end{ex}
\begin{ex}%[0D3N1-2] %5
	Trong các đẳng thức sau, đẳng thức nào đúng?
	\choice
	{ $\sin \left(180^\circ-\alpha\right) = \cos \alpha $}
	{ $\sin \left( 180^\circ -\alpha \right) =  - \cos \alpha $}
	{$\sin \left( 180^\circ-\alpha \right) =  - \sin \alpha $}
	{\True $\sin \left( 180^\circ - \alpha  \right) = \sin \alpha $}
	\loigiai{ Vì $\alpha$ và $180^\circ - \alpha$ là hai cung bù nhau nên theo tính chất giá trị lượng giác của hai cung bù nhau ta có đáp án.
	}
\end{ex}
\begin{ex}%[0D3N1-2] %6
	Cho  $\alpha$ là góc tù. Điều khẳng định nào sau đây là đúng?.
	\choice
	{ $\sin \alpha  < 0$}
	{ $\cos \alpha  > 0$}
	{\True $\tan \alpha  < 0$}
	{$\cot \alpha  > 0$}
	\loigiai{ Dựa vào giá trị lượng giá của cung tư $0^\circ$ đến $180^\circ$. Suy ra $\sin \alpha >0; \cos \alpha <0$ với $\alpha$ là góc tù, nên $\tan \alpha= \dfrac{\sin \alpha}{\cos \alpha} <0$.
	}
\end{ex}
\begin{ex}%[0D3N1-2] %7
	Cho hai góc nhọn $\alpha$ và $\beta$ trong đó  $\alpha < \beta$. Khẳng định nào sau đây là \textbf{sai}?
	\choice
	{ $\tan \alpha  + \tan \beta  > 0$}
	{\True $\cos \alpha  < \cos \beta $}
	{$\sin \alpha  < \sin \beta $}
	{Nếu $\alpha  + \beta  = 90^\circ$ thì $\cos \alpha  = \sin \beta $}
	\loigiai{
		$\alpha$ và $\beta$ là góc nhọn nên có điểm biểu diễn thuộc góc phần tư thứ nhất, có các giá trị lượng giác đều dương nên $\tan \alpha  + \tan \beta  > 0$, $\alpha  < \beta $ nên $\sin \alpha  < \sin \beta $, đáp án C đúng theo tính chất của hai góc phụ nhau
		
	}
\end{ex}
\begin{ex}%[0D3N1-2] %8
	Biết $A, B, C$ lần lượt là ba góc của tam giác $ABC$, mệnh đề nào sau đây là đúng?
	\choice
	{ ${\text{cos}}\,\left( {A + \,C} \right)\, = \,{\text{cos}}\,B$}
	{\True ${\text{tan}}\,\left( {A + \,C} \right)\, = \,\, - \tan \,B$}
	{${\text{cot}}\,\left( {A + \,C} \right)\, = \,{\text{cot}}\,B$}
	{${\text{sin}}\,\left( {A + \,C} \right)\, = \,\, - \sin \,B$}
	\loigiai{
		Xét phương án A: ${\text{cos}}\,\left( {A + \,C} \right)\, = \,\,{\text{cos}}\,\left( {\pi \, - \,\,B} \right)\, = \,\, - {\text{cos}}\,B\,\,\, \Rightarrow \,$ A sai\\
		Xét phương án B: ${\text{tan}}\,\left( {A + \,C} \right)\, = \,\,\tan \,\left( {\pi \, - \,\,B} \right)\, = \,\, - \tan \,B$  . Vậy B đúng\\
		Xét phương án C: ${\text{cot}}\,\left( {A + \,C} \right)\, = \,\,{\text{cot}}\,\left( {\pi \, - \,\,B} \right)\, = \,\, - {\text{cot}}\,B\,\,\, \Rightarrow \,$ C sai\\
		Xét phương án D: ${\text{sin}}\,\left( {A + \,C} \right)\, = \,\,\sin \,\left( {\pi \, - \,\,B} \right)\, = \,\,\sin \,B\,\,\, \Rightarrow \,$ D sai\\
	}
\end{ex}
\begin{ex}%[0D3N1-2] %9
	Biết $\cos \alpha =\dfrac{-2}{3}$. Tính $\tan \alpha$.
	\choice
	{ $\dfrac{5}{4}$}
	{ $ - \dfrac{5}{2}$}
	{$\dfrac{{\sqrt 5 }}{2}$}
	{\True $ - \dfrac{{\sqrt 5 }}{2}$}
	\loigiai{
		Do $\cos \alpha <0$ nên $\tan \alpha <0$.\\
		Ta có: $1 + {\tan ^2}\alpha  = \dfrac{1}{{{{\cos }^2}\alpha }} \Leftrightarrow {\tan ^2}\alpha  = \dfrac{5}{4} \Rightarrow \tan \alpha  =  - \dfrac{{\sqrt 5 }}{2}$. }
\end{ex}
\begin{ex}%[0D3N1-2] %10
	Cho $\cos x = \dfrac{1}{2}$. Tính giá trị của biểu thức $P = 3{\sin ^2}x + 4{\cos ^2}x$.
	\choice
	{  \True $\dfrac{{13}}{4}$}
	{ $\dfrac{7}{4}$}
	{$\dfrac{{11}}{4}$}
	{$\dfrac{{15}}{4}$}
	\loigiai{
		Ta có $P = 3{\sin ^2}x + 4{\cos ^2}x = 3\left( {{{\sin }^2}x + {{\cos }^2}x} \right) + {\cos ^2}x = 3 + {\left( {\frac{1}{2}} \right)^2} = \dfrac{{13}}{4}$  }
\end{ex}
\begin{ex}%[0D3N1-2] %11
	Cho biết $\cos \alpha = \dfrac{-2}{3}$. Tính giá trị của biểu thức $E= \dfrac{\cot \alpha +3 \tan \alpha}{2\cot \alpha +\tan \alpha}.$
	\choice
	{$ - \dfrac{{19}}{{13}}$}
	{\True $\dfrac{{19}}{{13}}$}
	{$\dfrac{{25}}{{13}}$}
	{$- \dfrac{{19}}{{13}}$}
	\loigiai{
		$E = \dfrac{{\cot \alpha  + 3\tan \alpha }}{{2\cot \alpha  + \tan \alpha }} = \dfrac{{1 + 3{{\tan }^2}\alpha }}{{2 + {{\tan }^2}\alpha }} = \dfrac{{3\left( {{{\tan }^2}\alpha  + 1} \right) - 2}}{{1 + \left( {1 + {{\tan }^2}\alpha } \right)}} = \dfrac{{\dfrac{3}{{{{\cos }^2}\alpha }} - 2}}{{\dfrac{1}{{{{\cos }^2}\alpha }} + 1}} = \dfrac{{3 - 2{{\cos }^2}\alpha }}{{1 + {{\cos }^2}\alpha }} = \dfrac{{19}}{{13}}$ }
\end{ex}
\begin{ex}%[0D3N1-2] %12
	Rút gọn biểu thức sau $A = \dfrac{{{{\cot }^2}x - {{\cos }^2}x}}{{{{\cot }^2}x}} + \dfrac{{\sin x.\cos x}}{{\cot x}}$.
	\choice
	{ $A = 4$}
	{ $A = 2$}
	{\True $A = 1$}
	{$A = 3$}
	\loigiai{
		$A = \dfrac{{{{\cot }^2}x - {{\cos }^2}x}}{{{{\cot }^2}x}} + \dfrac{{\sin x.\cos x}}{{\cot x}} = \dfrac{{\dfrac{{{{\cos }^2}x}}{{{{\sin }^2}x}} - {{\cos }^2}x}}{{\dfrac{{{{\cos }^2}x}}{{{{\sin }^2}x}}}} + \dfrac{{\sin x.\cos x}}{{\dfrac{{\cos x}}{{\sin x}}}} $\\
		$= \dfrac{{{{\cos }^2}x\left( {1 - {{\sin }^2}x} \right)}}{{{{\cos }^2}x}} + {\sin ^2}x = 1 - {\sin ^2}x + {\sin ^2}x = 1$
	}
\end{ex}
\Closesolutionfile{ans}
\indapan{6}{ans/ans-0-B1}
\cauds
\Opensolutionfile{ans}[ans/ans-0-B1-DS]
\begin{ex}%[0D3H1-5]
	\immini{
		Các mệnh đề sau \textbf{đúng} hay \textbf{sai}?
		\choiceTF
		{ $D = a\sin {90^\circ } + b\cos {90^\circ } + c\sin {180^\circ } = 2a$}
		{ $E = {\sin ^2}{60^\circ } + 2{\cos ^2}{30^\circ } - 5{\tan ^2}{45^\circ } = \dfrac{{11}}{4}$}
		{$F = {\cos ^2}{24^\circ } + {\cos ^2}{66^\circ } + {\cos ^2}{1^0} + {\cos ^2}{89^\circ } = 3$}
		{\True $G = {\cos ^2}{45^\circ } - 2{\cos ^2}{50^\circ } + 2{\sin ^2}{45^\circ } - 2{\cos ^2}{40^\circ } + 5\tan {55^\circ }\cot {125^\circ } = \dfrac{{ - 11}}{2}$}
	}{
		%\begin{tikzpicture}[scale=0.7,line join=round,line cap=round,font=\footnotesize,>=stealth]
		%\draw[step=1,gray,very thin] (-3.5,-1.5) grid (5,5);
		%	\draw[->] (-3.5,0)--(5,0); %Trục hoành
		%	\draw[->] (0,-1.5)--(0,5); %Trục tung
		%	\draw (4.7,0) node[above]{$x$} (0,4.7) node[left]{$y$} (0,0) node[below left]{$O$}; %Tên trục
		%	\draw[fill=black] (0,0) circle(0.03);
		%	\foreach \x in {3}{
			%	\draw[fill=black] (\x,0) node[below]{$\x$} circle(0.01);
			%	};%Đánh số trục hoành
		%	\foreach \y in {-1}{
			%		\draw[fill=black] (0,\y) node[right]{$\y$} circle(0.01);
			%	};%Đánh số trục tung
		%	\fill (-3,0) node[above left]{$-3$} circle(1 pt)
		%	 (-2,0) node[above left]{$-2$} circle(1pt)
		%	 (0,1) node[below right]{$1$} circle(1pt)
		%	 (0,4) node[left]{$4$} circle(1pt);
		%	\draw[dashed] (-3,0)--(-3,-1)--(0,-1); 
		%	\draw[dashed] (3,0)--(3,4)--(0,4);
		%	\draw (-3,-1)--(-1,1)--(0,1)--(3,4);			
		%\end{tikzpicture}
	}
	\loigiai{
		\begin{itemchoice}
			\itemch Sai. $D = a.1 + b.0 + c.0 = a$.
			\itemch Sai. $E = {\left( {\dfrac{{\sqrt 3 }}{2}} \right)^2} + 2 \cdot {\left( {\dfrac{{\sqrt 3 }}{2}} \right)^2} - 5.1 =  - \dfrac{{11}}{4}$
			\itemch Sai. $F = {\cos ^2}{24^0} + {\sin ^2}{24^0} + {\cos ^2}{1^0} + {\sin ^2}{1^0} = 1 + 1 = 2$.
			\itemch Đúng. $G = {\cos ^2}{45^\circ } + {\sin ^2}{45^\circ } + {\sin ^2}{45^\circ } - 2\left( {{{\cos }^2}{{50}^\circ } + {{\cos }^2}{{40}^\circ }} \right) - 5\tan {55^\circ }\cot {55^\circ }$\\
			$ = {\cos ^2}{45^\circ } + {\sin ^2}{45^\circ } + {\sin ^2}{45^\circ } - 2\left( {{{\cos }^2}{{50}^\circ } + {{\sin }^2}{{50}^\circ }} \right) - 5\tan {55^\circ }\cot {55^\circ } = \dfrac{{ - 11}}{2}$.
		\end{itemchoice}
	}
\end{ex}
\begin{ex}%[0D3H1-5] % câu 2
		Cho $\cos \alpha  =  - \dfrac{3}{4}\left( {{0^\circ } < \alpha  < {{90}^\circ }} \right)$. Khi đó:
		\choiceTF
		{\True ${\sin ^2}\alpha  = \dfrac{7}{{16}}$}
		{ $\sin \alpha  < 0$}
		{$\sin \alpha  =  - \dfrac{{\sqrt 7 }}{4}$}
		{\True $\cot \alpha  =  - \dfrac{{3\sqrt 7 }}{7}$}
	\loigiai{
		Vì ${\sin ^2}\alpha  + {\cos ^2}\alpha  = 1 \Rightarrow {\sin ^2}\alpha  = 1 - {\cos ^2}\alpha  = 1 - {\left( { - \dfrac{3}{4}} \right)^2} = \dfrac{7}{{16}}$\\
		Mà ${0^\circ } < \alpha  < {90^\circ }$ nên $\sin \alpha >0 $. Do đó $\sin \alpha  = \sqrt {\dfrac{7}{{16}}}  = \dfrac{{\sqrt 7 }}{4}$\\
		$\tan \alpha  = \dfrac{{\sin \alpha }}{{\cos \alpha }} = \dfrac{{\dfrac{{\sqrt 7 }}{4}}}{{ - \dfrac{3}{4}}} =  - \dfrac{{\sqrt 7 }}{3};\cot \alpha  = \dfrac{{\cos \alpha }}{{\sin \alpha }} =  - \dfrac{{3\sqrt 7 }}{7}$
	}
\end{ex}
\begin{ex}%[0D3H1-5] % câu 3
		Cho góc $\alpha$ thỏa mãn $\sin \alpha = \dfrac{3}{5}.$ Khi đó
		\choiceTF
		{\True ${\sin ^2}\alpha  = \dfrac{9}{{25}}$}
		{\True ${\cos ^2}\alpha  = \dfrac{{16}}{{25}}$}
		{\True $\dfrac{{\cot \alpha  + \tan \alpha }}{{\cot \alpha  - \tan \alpha }} = \dfrac{{25}}{7}$}
		{$\dfrac{1}{{{{\cos }^2}\alpha  - {{\sin }^2}\alpha }} = \dfrac{7}{{25}}$}
	\loigiai{
		Ta có $H = \dfrac{{\cot \alpha  + \tan \alpha }}{{\cot \alpha  - \tan \alpha }} = \dfrac{{\dfrac{{\cos \alpha }}{{\sin \alpha }} + \dfrac{{\sin \alpha }}{{\cos \alpha }}}}{{\dfrac{{\cos \alpha }}{{\sin \alpha }} - \dfrac{{\sin \alpha }}{{\cos \alpha }}}} = \dfrac{{\dfrac{{{{\cos }^2}\alpha  + {{\sin }^2}\alpha }}{{\sin \alpha  \cdot \cos \alpha }}}}{{\dfrac{{{{\cos }^2}\alpha  - {{\sin }^2}\alpha }}{{\sin \alpha  \cdot \cos \alpha }}}} = \dfrac{1}{{{{\cos }^2}\alpha  - {{\sin }^2}\alpha }}.$\\
		Vì $\sin \alpha = \dfrac{3}{5}$ nên $\sin^2 \alpha = \dfrac{9}{25}$ và ${\cos ^2}\alpha  = 1 - {\sin ^2}\alpha  = \dfrac{{16}}{{25}}$. 
		
		Do đó: $H = \dfrac{1}{{{{\cos }^2}\alpha  - {{\sin }^2}\alpha }} = \dfrac{1}{{\dfrac{{16}}{{25}} - \dfrac{9}{{25}}}} = \dfrac{{25}}{7}$
	}
\end{ex}
\begin{ex}%[0D3H1-5] % câu 4
		Cho $\cot \alpha = - \sqrt 2 , (0^\circ < \alpha < 180^\circ) $. Khi đó: 
		\choiceTF
		{\True $\sin \alpha >0$}
		{$\sin \alpha  =  \pm \dfrac{1}{{\sqrt 3 }}$}
		{\True $\cos \alpha = - \dfrac{\sqrt 6}{3}$}
		{$\tan \alpha = \dfrac{1}{\sqrt 2}$}
	\loigiai{
		Ta có: ${\sin ^2}\alpha  = \dfrac{1}{{1 + {{\cot }^2}\alpha }} = \dfrac{1}{{1 + 2}} = \dfrac{1}{3} \Leftrightarrow \sin \alpha  =  \pm \dfrac{1}{{\sqrt 3 }}$
		Do $0^\circ < \alpha < 180^\circ$ nên $\sin \alpha >0 \Rightarrow \sin \alpha = \dfrac{1}{\sqrt 3}$.\\
		Mà $\cot \alpha = \dfrac{\cos \alpha}{\sin \alpha} \Rightarrow \cos \alpha = \cot \alpha . \sin \alpha = - \sqrt 2 . \dfrac{1}{\sqrt 3}=-\dfrac{\sqrt 6}{3}$
	}
\end{ex}
\begin{ex}%[0D3H1-5] % câu 5
		Xét tính \textbf{đúng}, \textbf{sai} của các đẳng thức sau:
		\choiceTF
		{\True ${\left( {{{\sin }^2}\alpha  + {{\cos }^2}\alpha } \right)^2} = 1$}
		{\True $1 + 2\sin \alpha  \cdot \sin \alpha  = {(\sin \alpha  + \cos \alpha )^2}$}
		{\True $1 - 2\sin \alpha  \cdot \sin \alpha  = {(\sin \alpha  - \cos \alpha )^2}$}
		{\True $1 - 2{\sin ^2}\alpha  \cdot {\sin ^2}\alpha  = \sin {\alpha ^4} + {\cos ^4}\alpha $}
	\loigiai{
		\begin{itemchoice}
			\itemch ${\left( {{{\sin }^2}\alpha  + {{\cos }^2}\alpha } \right)^2} = 1$
			\itemch $1 + 2\sin \alpha  \cdot \cos \alpha  = {\sin ^2}\alpha  + {\cos ^2}\alpha  + 2\sin \alpha  \cdot \cos \alpha  = {(\sin \alpha  + \cos \alpha )^2}$
			\itemch $1 - 2\sin \alpha  \cdot \cos \alpha  = {\sin ^2}\alpha  + {\cos ^2}\alpha  - 2\sin \alpha  \cdot \cos \alpha  = {(\sin \alpha  - \cos \alpha )^2}$
			\itemch Ta có: ${\sin ^4}\alpha  + {\cos ^4}\alpha  + 2{\sin ^2}\alpha  \cdot {\cos ^2}\alpha  = {\left( {{{\sin }^2}\alpha  + {{\cos }^2}\alpha } \right)^2} = 1$\\
			Do đó: $1 - 2{\sin ^2}\alpha {\cos ^2}\alpha  = {\sin ^4}\alpha  + {\cos ^4}\alpha $.
		\end{itemchoice}
	}
\end{ex}
\begin{ex}%[0D3H1-5] % câu 6
		Xét tính \textbf{đúng}, \textbf{sai} của các đẳng thức sau:
		\choiceTF
		{\True $A = 4\sin {30^\circ } + {(\sqrt 3 )^3} \cdot \tan {30^\circ } = 5$}
		{\True $B = \dfrac{1}{{\sqrt 3 }}\cos {30^\circ } - 3\sqrt 2 \sin {45^\circ } + \cot {45^\circ } = \dfrac{3}{2}$}
		{ $C = {\sin ^2}{60^\circ } + {\tan ^2}{30^\circ } - 2 =  - \dfrac{{11}}{{12}}$}
		{\True $D = 2\left( {\sin {{150}^\circ } + \sqrt 3 } \right) + \dfrac{1}{{2\sqrt 2 }}\cos {135^\circ } - 3\tan {150^\circ }  = \dfrac{3}{4} + 2\sqrt 3 $}	
	\loigiai{
		\begin{itemchoice}
			\itemch $A = 4 \cdot \dfrac{1}{2} + 3\sqrt 3  \cdot \dfrac{1}{{\sqrt 3 }} = 2 + 3 = 5$
			\itemch $B = \dfrac{1}{{\sqrt 3 }} \cdot \dfrac{{\sqrt 3 }}{2} - 3\sqrt 2  \cdot \dfrac{{\sqrt 2 }}{2} + 1 = \dfrac{1}{2} - 3 + 1 =  - \dfrac{3}{2}$
			\itemch $C = {\left( {\dfrac{{\sqrt 3 }}{2}} \right)^2} + {\left( {\dfrac{1}{{\sqrt 3 }}} \right)^2} - 2 = \dfrac{3}{4} + \dfrac{1}{3} - 2 =  - \dfrac{{11}}{{12}}$
			\itemch $D = 2\sin {30^\circ } + 2\sqrt 3  - \dfrac{1}{{2\sqrt 2 }}\cos {45^\circ } + 3\tan {30^\circ }$ do:\\
			$\left\{ {\begin{array}{*{20}{l}}
					\begin{gathered}
						\sin {150^\circ } = \sin \left( {{{180}^\circ } - {{30}^\circ }} \right) = \sin {30^\circ } \hfill \\
						\cos {135^\circ } = \cos \left( {{{180}^\circ } - {{45}^\circ }} \right) =  - \cos {45^\circ } \hfill \\
						\tan {150^\circ } = \tan \left( {{{180}^\circ } - {{30}^\circ }} \right) =  - \tan {30^\circ } \hfill \\ 
					\end{gathered}  
			\end{array}} \right.$\\
			Do đó: $D = 2 \cdot \dfrac{1}{2} + 2\sqrt 3  - \dfrac{1}{{2\sqrt 2 }} \cdot \dfrac{{\sqrt 2 }}{2} + 3 \cdot \dfrac{{\sqrt 3 }}{3} = 1 + 2\sqrt 3  - \dfrac{1}{4} + \sqrt 3  = \dfrac{3}{4} + 3\sqrt 3 .$
		\end{itemchoice}
	}
\end{ex}
\Closesolutionfile{ans}
\indapan{3}{ans/ans-0-B1-DS}

\Opensolutionfile{ans}[ans/ans-0-B1-KQ]
\caukq
\begin{ex}%[0D3N1-3] %cau 1
	Tính giá trị của biểu thức sau:\\
	$A = {\cos ^2}{15^\circ } + {\cos ^2}{25^\circ } + {\cos ^2}{35^\circ } + {\cos ^2}{45^\circ } + {\sin ^2}{15^\circ } + {\sin ^2}{25^\circ } + {\sin ^2}{35^\circ }$.
	\shortans{$3{,}5$}
	\loigiai{
		$A = \left( {{{\cos }^2}{{15}^\circ } + {{\sin }^2}{{15}^\circ }} \right) + \left( {{{\cos }^2}{{25}^\circ } + {{\sin }^2}{{25}^\circ }} \right) + \left( {{{\cos }^2}{{35}^\circ } + {{\sin }^2}{{35}^\circ }} \right) + {\cos ^2}{45^\circ }$\\
		$ = 1 + 1 + 1 + {\left( {\dfrac{{\sqrt 2 }}{2}} \right)^2} = \dfrac{7}{2}$
	}
\end{ex}
\begin{ex}%[0D3N1-3] %cau 2
	Tính giá trị cảu biểu thức sau:\\
	$B = \tan {1^\circ }\tan {2^\circ }\tan {3^\circ } \ldots \tan {89^\circ }$.
	\shortans{$1$}
	\loigiai{
		$B = \tan {1^\circ }\tan {2^\circ }\tan {3^\circ } \ldots \tan {89^\circ }$\\
		$ = \left( {\tan {1^\circ }\tan {{89}^\circ }} \right) \cdot \left( {\tan {2^\circ }\tan {{88}^\circ }} \right) \ldots  \ldots \left( {\tan {{44}^\circ }\tan {{46}^\circ }} \right) \cdot \tan {45^\circ }$\\
		$ = \left( {\tan {1^\circ } \cdot \cot {1^\circ }} \right) \cdot \left( {\tan {2^\circ }\cot {2^\circ }} \right) \ldots  \ldots \left( {\tan {{44}^\circ }\cot {{44}^\circ }} \right) \cdot \tan {45^\circ } = 1.1 \ldots  \ldots 1 = 1$.
	}
\end{ex}
\begin{ex}%[0D3N1-3] %cau 3
	
	Cho $\cot \alpha = \dfrac{1}{3}.$ Tính giá trị của biểu thức $A = \dfrac{{3\sin \alpha  + 4\cos \alpha }}{{2\sin \alpha  - 5\cos \alpha }}$.
	\shortans{$13$}
	\loigiai{
		Do $\cot \alpha  = \dfrac{{\cos \alpha }}{{\sin \alpha }} = \dfrac{1}{3} \Rightarrow \sin \alpha  \ne 0$\\
		Chia hai vế của biểu thức $A$ cho $\sin \alpha$, ta có: 
		
		
		$A = \dfrac{{3 + 4\dfrac{{\cos \alpha }}{{\sin \alpha }}}}{{2 - 5\dfrac{{\cos \alpha }}{{\sin \alpha }}}} = \dfrac{{3 + 4\cot \alpha }}{{2 - 5\cot \alpha }} = \dfrac{{3 + 4 \cdot \dfrac{1}{3}}}{{2 - 5 \cdot \dfrac{1}{3}}} = 13 \, \hfill$
		
		
		
	}
\end{ex}
\begin{ex}%[0D3N1-3] %cau 4
	
	Cho $\sin x +\cos x =m$. Tính theo $m$ giá trị của biểu thức $M=\sin x . \cos x$ ta thu được kết quả dạng $M = \dfrac{m^2-a}{b}$ với $a,b\in \mathbb{N}$. Tính $a+b$.
	\shortans{$3$}
	\loigiai{
		$\begin{gathered}
			\sin x + \cos x = m\,\, \Rightarrow \,{\left( {\sin x + \cos x} \right)^2}\, = {m^2} \hfill \\
			\, \Leftrightarrow \,\left( {{{\sin }^2}x + {{\cos }^2}x} \right) + 2\sin x.\cos x = {m^2}\, \hfill \\ 
			\Leftrightarrow \,1 + 2\sin x.\cos x = {m^2}\, \hfill \\
			\Leftrightarrow \,\sin x.\cos x = \dfrac{{{m^2} - 1}}{2} \, \hfill \\
			\text{Vậy } M=\dfrac{m^2 -1}{2} \, \hfill
		\end{gathered} $.
	Vậy $a=1,b=2$ nên $a+b=3$.
	}
\end{ex}
\begin{ex}%[0D3N1-3] %cau 5
	Cho $\tan \alpha +\cot \alpha =m$. Biết $m_0$ là giá trị để $\tan ^2 \alpha + \cot ^2 \alpha =7$. Tính $m_0^2+2015$.
	\shortans{$2024$}
	\loigiai{
		$7 = {\tan ^2}\alpha  + {\cot ^2}\alpha  = {\left( {\tan \alpha  + \cot \alpha } \right)^2} - 2 \Rightarrow {m^2} = 9 \Leftrightarrow m =  \pm 3$. Vậy $m_0^2+2015=2024$.
	}
\end{ex}
\begin{ex}%[0D3N1-3] %cau 6
	Giá trị của biểu thức $B = \dfrac{{2{{\sin }^2}{{30}^\circ }}}{{1 - 2{{\cos }^2}{{30}^\circ }}} + 4\sin {60^\circ } - \cot {30^\circ }$ có dạng $a+b\sqrt{3}$ trong đó $a,b\in \mathbb{Z}$. Tính $a+b$.
	\shortans{$0$}
	\loigiai{
		$B = \dfrac{{2 \cdot {{\left( {\dfrac{1}{2}} \right)}^2}}}{{1 - 2{{\left( {\dfrac{{\sqrt 3 }}{2}} \right)}^2}}} + 4 \cdot \dfrac{{\sqrt 3 }}{2} - \sqrt 3  =  - 1 + \sqrt 3 $. Vậy $a=-1; b=1$ nên $a+b=0$.
	}
\end{ex}

\Closesolutionfile{ans}
\indapan{6}{ans/ans-0-B1-KQ}
